% !Mode:: "TeX:UTF-8"
\chapter{绪论}
\label{cha:intro}

\section{课题来源}
\label{sec:source}
本课题来源于国家自然科学基金优秀青年基金项目 ``复杂非线性系统的模糊控制方法与应
用'' (项目编号: 62122046), 以及该项目资助下开展的 ``无人船协同与对抗'' 研究方向. 本
研究旨在针对复杂海洋环境中欠驱动水面无人艇的博弈策略问题, 开展基于深度强化学习的
决策方法研究, 并完成多无人艇物理实验平台的搭建与实物实验验证.

\section{研究背景及意义}
\label{sec:background}
水面无人艇 (Unmanned Surface Vehicle, USV) 是一种具备遥控或自主航行能力的水面运
载平台, 具有成本低、机动灵活、部署快捷、无人员伤亡风险等显著优势, 是现代海洋科技
装备智能化、无人化的重要载体~\cite{peng2021usv,liao2014usv}. 随着传感器、通信、计
算与人工智能技术的迅猛发展, 无人艇已广泛应用于海洋环境监测、水质巡查、航道测量、
海上搜救、水面清洁等民用场景, 以及巡逻警戒、电子对抗、反潜扫雷、目标拦截等军事任
务~\cite{jin2018marine,roberts2008marine}. 作为衡量一个国家海洋科技实力的重要标
志, 无人艇技术正日益成为国际海洋装备竞争的前沿方向.

近年来, 随着海战形态向智能化、体系化、无人化方向加速演进, 传统单艇遥控或预规划航
迹已难以满足日益复杂的海上对抗任务需求. 多无人艇协同与对抗成为新一代海洋无人装备
的关键技术~\cite{peng2021usv}. 在多无人艇执行协同博弈任务时, 智能体既需要考虑队
形保持、协同避障等传统编队问题, 又必须根据对手的运动状态、己方的任务能力以及敌我
双方的相对态势, 实时作出策略性决策, 以完成诸如追击逃逸目标、合围围捕敌方、护航重
要设施等典型博弈任务~\cite{su2022confrontation,xia2023cooperativeusv}. 因此, 围绕
复杂环境下欠驱动无人艇的博弈策略研究, 不仅是多智能体协同控制的前沿课题, 也是服务
于海洋强国战略、支撑智能化海上作战体系建设的重要基础.

\begin{figure}[!htbp]
  \centering
  % \includegraphics[scale=0.4]{usv_overview.pdf}
  \vspace{2.5cm}
  \bicaption{典型水面无人艇平台示意}{Illustration of typical USV platforms}
  \label{fig:usv_overview}
\end{figure}

然而, 复杂海洋环境下无人艇博弈策略的研究面临诸多理论与技术挑战, 可概括为以下三个
层面. 第一, 从\textbf{载体层面}看, 出于结构简化与成本控制的考虑, 绝大多数水面无人艇
仅配备螺旋桨推进器与方向舵, 即仅有纵向推力和艏向力矩两个控制输入, 呈现典型的欠驱
动动力学特征; 同时, 无人艇的六自由度运动存在高度非线性耦合, 传统线性或解耦控制方
法难以获得理想的性能. 第二, 从\textbf{环境层面}看, 海洋环境中的风、浪、流等外部扰
动强度大、时变性显著, 难以被精确建模; 海面上的暗礁、浮标、他船等障碍物分布具有随
机性和动态性; 水声、无线电等通信链路易受环境干扰, 信息传输存在时延和丢包. 第三, 从
\textbf{博弈层面}看, 对手的策略具有高度不确定性与对抗性, 多智能体间策略相互耦合导
致学习过程呈现非平稳特性, 训练收敛慢、策略泛化性差等问题突出~\cite{drl_2018}. 上
述三个层面的挑战彼此交织, 使得传统基于精确建模与最优控制的方法难以直接应用, 亟需
探索能够在欠驱动约束、环境不确定性和多智能体非平稳性下仍能高效决策的策略优化方法.

近年来, 深度强化学习 (Deep Reinforcement Learning, DRL) 以其无需精确系统模型、能
够通过与环境交互自主学习最优策略的特点, 为复杂博弈决策问题提供了崭新的解决思路
~\cite{sutton1998introduction,drl_2018}. 将深度强化学习引入水面无人艇博弈策略研
究, 有望在欠驱动、强扰动、对抗性环境下实现高效、鲁棒的自主决策. 因此, 开展复杂环境
下基于深度强化学习的无人艇博弈策略研究, 对于推进海上无人装备智能化水平、支撑国家
海洋安全战略具有重要的理论价值和工程意义.

\section{国内外研究现状}
\label{sec:review}

\subsection{水面无人艇智能控制研究现状}
\label{subsec:usv_ctrl}
水面无人艇的研究最早可追溯至 1993 年美国麻省理工学院开展的 ``ARTEMIS'' 项目, 该项
目建造了能够自主航行的小型科学考察船, 开启了现代水面无人艇研究的序幕. 进入 21 世
纪后, 随着微机电系统、嵌入式计算与无线通信等技术的成熟, 无人艇的研究重心从单艇自
主航行逐步拓展到多艇协同感知、协同决策与协同作战~\cite{peng2021usv,jin2018marine}.

在国外, 美国麻省理工学院研制的 ``SCOUT'' 水面无人艇搭载了 MOOS-IvP 自主软件与采样
平台, 成功完成了港口监控、水文测量等多无人艇协同任务. 美国海军研究署
(Office of Naval Research, ONR) 曾开展过无人艇蜂群战术测试, 其研制的 ``CARACaS''
分布式智能指挥系统使四艘无人艇能够在无人工干预条件下完成区域搜索、识别、跟踪与常
规巡逻, 标志着美军已初步具备多无人艇集群协同作战能力. 葡萄牙里斯本大学 IST-ISR 研
究所围绕协同异构无人系统项目, 对多无人艇协同避障、软硬件联合运动控制、时变通信拓
扑等关键技术开展了系统研究, 并在大西洋沿岸开展了多次海上实验. 此外, 法国、英国、以
色列等国也在加速无人艇集群相关的研发与海试, 推动该技术从实验室走向实战.

在国内, 以华中科技大学研制的 ``HUSTER-68'' 无人艇、大连海事大学的 ``蓝信号-I'' 无人
艇、哈尔滨工程大学的 ``海豚'' 系列无人艇为代表, 已经成功开展了多艇编队、协同巡航
等典型任务. 在产业界, 珠海云洲智能科技有限公司于 2018 年组织了 81 艘无人艇横穿港
珠澳大桥的 ``启航'' 编队表演, 并在万山海域完成了 56 艘无人艇参与的集群编队、协同避
障、容错控制等大规模测试, 展示了国产无人艇集群技术的实力. 尽管如此, 与无人机集群、
多机器人集群等领域相比, 我国无人艇集群技术的发展相对滞后. 其主要原因在于: 一是无人
艇欠驱动特性使控制器设计复杂度显著提高; 二是海面风浪、洋流等环境因素多变且难以建
模; 三是水面通信易受干扰、带宽有限; 四是集群系统自身的高维性与协同复杂性使算法设
计面临维度灾难.

从发展趋势上看, 国内外水面无人艇研究正呈现以下三个特点: (i) 由 ``单艇自主'' 向
``多艇协同'' 发展, 集群规模逐步扩大; (ii) 由 ``机械式编队'' 向 ``智能化博弈对抗'' 演
进, 强调对抗性环境下的策略决策能力; (iii) 由 ``固定任务'' 向 ``跨域协同''
扩展, 与无人机、无人潜航器等平台联合作业. 这些趋势共同呼唤更加智能、鲁棒、可迁
移的无人艇博弈策略算法.

\subsection{无人艇博弈策略研究现状}
\label{subsec:game_review}
无人艇博弈问题本质上是在连续状态-连续动作空间下多智能体交互的动态决策问题. 围绕
此类问题, 国内外学者提出了多种研究范式, 本节按照主流的分类方式归纳为确定性方法、
几何方法、启发式方法、基于力的方法与强化学习方法五大类, 并逐一分析其代表性工作与
局限性.

\subsubsection{确定性方法}
确定性方法起源于 Isaacs 在 20 世纪中期提出的微分博弈理论, 其核心思想是将博弈问题
转化为求解 Hamilton-Jacobi-Isaacs (HJI) 偏微分方程, 从而获得纳什均衡意义下的最优
策略~\cite{isaacs1999differential}. 在多对多追逃博弈方面,
Lopez 等人~\cite{lopez2019multiagent} 通过图论对多智能体间通信拓扑进行建模,
提出了能够在有限时间内捕获逃逸者的协同策略, 并分析了策略的渐近收敛性质.
Salimi 等人~\cite{salimi2019differential} 针对多个追击者对单一逃逸者的微分博弈问题,
设计了基于追逐距离的回报函数, 利用最优控制理论推导出最优追击策略.
Garcia 等人~\cite{garcia2020multi} 进一步提出了一种连续可微且满足 HJI 方程的价值
函数, 一定程度上缓解了传统 HJI 方法面临的维度灾难问题, 并针对数量占优的追击者设计
了协同制导律. 整体而言, 确定性方法具有严格的数学理论基础和可证明的最优性, 然而其
求解过程对状态空间维度高度敏感, 难以直接应用于包含多个追击者、多个逃逸者以及复杂
障碍物的高维博弈场景.

\subsubsection{几何方法}
几何方法通过构造描述智能体相对位置与可达区域的几何结构 (如 Voronoi 图、阿波罗尼
奥斯圆、缓冲 Voronoi 单元等) 来指导策略决策, 具有直观性强、计算量小的特点.
Von Moll 等人~\cite{vonmoll2019multi} 基于 Voronoi 图提出了多对一追逃博弈的几何
解法, 为逃逸者推导出了一种不被捕获的开环最优策略.
Ramana 等人~\cite{ramana2017pursuit} 利用阿波罗尼奥斯圆构建包围模型, 设计了针对
高速逃逸者的逃逸策略, 揭示了逃逸者可通过编队间隙进行机动逃逸的条件.
Tian 等人~\cite{tian2021distributed} 研究了含障碍环境下的分布式追捕问题, 提出了
基于 Voronoi 单元内最近点搜索的追击与避障策略. 几何方法的优势在于能够将高维博弈
问题降维到几何约束的求解上, 但其对智能体动力学约束、外部扰动等因素的表达能力较弱,
难以应用于真实海洋环境中的欠驱动无人艇博弈.

\subsubsection{启发式方法}
启发式方法通过模拟自然界生物群体的捕食或躲避行为, 设计具有生物学可解释性的博弈策
略. Ni 等人~\cite{ni2011bioinspired} 最早提出了基于生物启发神经网络的多机器人协
同围捕方法, 将周围环境建模为神经活动场, 利用神经元活性引导智能体运动. 在此基础
上, 后续研究进一步将生物启发模型扩展至洋流影响下的水下机器人路径规划与多自主水
下航行器 (AUV) 任务分配. 此外, 受鹰类捕猎行为启发, 一些研究提出了分组攻击策略; 受
粒子群智能启发, MahmoudZadeh 等人提出了用于海洋采样任务的多无人艇不间断无碰撞路
径规划系统. 然而, 启发式方法普遍存在以下不足: 在复杂障碍场景下效率较低, 无法保证
策略的最优性, 对初始条件依赖性强, 且可扩展性有限.

\subsubsection{基于力的方法}
基于力的方法借鉴人工势场思想, 通过为智能体施加不同种类的 ``力'' 来引导其实时运动.
Angelani 利用群体互动模型分析集体追逐问题, 通过模拟群体成员之间的相互作用提升追
逐成功率. Janosov 等人提出了局部交互规则, 以简单协调策略优化多个追逐者的配合.
Fang 等人~\cite{fang2022cooperative} 引入了基于邻居距离的潜在力, 使追逐者能够形
成包围圈并逐步接近躲避者, 实现了对速度更快的自由运动逃逸者的协同追捕. 基于力的方
法结构简单、实时性好, 然而其本质上属于贪婪式决策, 缺乏对长期回报的显式优化, 且较
难处理智能体之间的深度合作与全局协同问题.

\subsubsection{强化学习方法}
强化学习通过智能体与环境的反复交互来学习最优策略, 不依赖精确的系统模型, 能够适应
动态变化的博弈场景, 近年来已成为无人艇博弈研究中最活跃的方向~\cite{drl_2018}. 在
追逃微分博弈中, Wang 等人~\cite{wangy2020cooperative} 使用深度确定性策略梯度算
法同时训练追击者与逃逸者的策略, 避免了传统方法中假设逃逸者采用固定策略的局限性.
在协同围捕方面, Lowe 等人~\cite{lowe2017maddpg} 提出的多智能体深度确定性策略梯
度 (MADDPG) 算法通过集中训练-分布式执行框架, 为多智能体博弈提供了通用解决方案;
在此基础上, 一些研究结合通信机制设计了分布式协同围捕策略.
Xia 等人~\cite{xia2023cooperative} 将 MADDPG 算法应用于多无人艇协同围捕场景,
针对连续动作空间、无边界环境和部分可观测性等实际问题进行了针对性改进.
Selvakumar 和 Bakolas~\cite{selvakumar2022minmax} 提出了 min-max Q-learning 方
法, 将高维状态空间映射至低维流形, 有效处理了多玩家追逃博弈问题. 此外, 近期兴起的
多智能体近端策略优化 (MAPPO) 算法由于其训练稳定性与收敛效率, 已被证实在多种多智
能体协作场景下具有出色的性能~\cite{yu2022surprising}.

尽管强化学习方法在无人艇博弈研究中取得了一系列进展, 但仍面临以下突出问题: (i) 大多
数现有工作忽视了海洋环境扰动在观测空间中的显式建模, 对环境不确定性的鲁棒性有限;
(ii) 针对协同围捕等复杂博弈任务, 现有方法普遍缺乏严格的成功判据, 导致奖励设计与算
法评价依据模糊, 训练过程容易陷入稀疏奖励困境; (iii) 现有研究多止步于数值仿真, 面向
真实水面环境的物理平台验证十分稀缺, 算法的工程可实现性有待考察.

\subsection{多智能体深度强化学习研究现状}
\label{subsec:marl_review}
多智能体深度强化学习 (Multi-Agent Deep Reinforcement Learning, MADRL) 是研究多
个智能体在共享环境中通过交互学习最优联合策略的前沿方向. 由于多智能体系统中各智能
体的策略相互耦合、环境从单个智能体视角呈现出显著的非平稳性, MADRL 算法的设计与
训练面临比单智能体强化学习更大的挑战.

在算法范式上, MADRL 主要遵循三种框架: 完全集中式 (Centralized)、完全分布式
(Decentralized) 和集中训练-分布式执行 (Centralized Training with Decentralized
Execution, CTDE). 其中, CTDE 范式由于兼顾训练稳定性与执行效率, 已成为当前主流.
Lowe 等人提出的 MADDPG 算法~\cite{lowe2017maddpg} 是 CTDE 范式的经典代表, 其
为每个智能体配备独立的 Actor 和 Critic, 在训练阶段 Critic 利用全局状态信息以克服
非平稳性, 在执行阶段 Actor 仅依赖局部观测实现分布式决策.
Yu 等人~\cite{yu2022surprising} 将 PPO 算法扩展至多智能体场景, 提出的 MAPPO 算
法在多种协作任务上超越了 MADDPG 与 QMIX 等基线方法, 表明 on-policy 方法在多智能
体场景下具有独特优势.

为了进一步提升多智能体策略的可扩展性与表达能力, 近年来研究者将注意力机制
~\cite{vaswani2017attention} 引入 MADRL 框架. Iqbal 等人~\cite{iqbal2019maac}
提出的 MAAC 算法在 Critic 网络中引入注意力机制, 使智能体能够自适应地关注影响自身
决策的关键邻居. 这一思想启发了后续一系列工作, 包括基于图注意力的 G2ANet, 基于
Transformer 的多智能体结构等. 对于本文关注的无人艇协同围捕博弈, 引入多头注意力机
制尤为必要, 因为多艇之间、艇-障碍之间、艇-目标之间的相互作用存在多种不同层次的依
赖关系, 单一的注意力子空间难以充分表达.

此外, 在奖励稀疏、训练困难等问题上, 后见经验回放 (Hindsight Experience Replay,
HER)、分层强化学习、基于课程学习的训练策略等也取得了一系列进展. 对于涉及对抗与博
弈的多智能体问题, 如何设计兼顾探索效率与策略鲁棒性的训练机制, 仍是当前研究的热点.

\subsection{目标-进攻-防御博弈研究现状}
\label{subsec:tad_review}
目标-进攻-防御博弈 (Target-Attacker-Defender, TAD) 作为一类特殊的非对称博弈, 因
其在反舰导弹拦截、关键目标防卫、空天对抗等领域的广泛应用而受到越来越多的关注. TAD
博弈的核心特征在于三方交互: 进攻者试图捕获目标, 防御者则在保护目标的同时拦截进攻
者, 三者行为彼此耦合.

在确定性框架下, Garcia 等人~\cite{garcia2018cooperative,garcia2019target} 将 TAD
博弈建模为线性二次微分博弈, 基于 Riccati 方程求解最优反馈策略, 并针对具有速度优势
的进攻者推导了目标保护的临界条件. 一些研究进一步将 TAD 博弈分解为子博弈, 利用分
阶段最优控制方法求解具有不同约束类型的复合制导律. 然而, 上述工作大多局限于简化的
三方场景, 较少考虑多对多的复杂对抗、非自杀式防御约束及其在欠驱动动力学下的求解.

基于强化学习的 TAD 研究近年来初见端倪, 主要关注进攻者或防御者策略的学习. 部分研
究者采用改进的软演员-评论家算法, 结合渗透奖励机制与长短时记忆网络, 提高了进攻者
在防御环境中的突破能力. 另一部分工作则通过设计基于行为的奖励函数, 结合辅助修正机
制, 训练出具备轨道机动能力的进攻者. 然而, 面向非自杀式防御者的多对多 TAD 博弈策
略设计尚鲜有报道. 非自杀式约束要求防御者在完成拦截后仍保留机动能力以持续保护目
标, 这一约束的引入使得传统 ``拦截即结束'' 的建模假设不再成立, 对策略学习和奖励设计
提出了新的挑战. 此外, 现有工作鲜有利用进攻者未来状态的预测来增强防御者的前瞻性决
策能力, 这也是本文第五章重点探讨的问题.

\section{现有研究存在的问题}
\label{sec:gap}
综合上述国内外研究现状分析可以看出, 尽管复杂环境下无人艇博弈策略研究近年来取得了
一系列重要进展, 但在理论模型、算法设计和工程验证等方面仍存在诸多亟待解决的问题,
可概括如下.

\textbf{问题 1: 欠驱动无人艇在扰动环境下的追逃策略鲁棒性不足.} 现有基于强化学习的
无人艇追逃研究, 大多将观测空间局限于智能体的运动状态和目标位置, 而忽视了海洋环境
中风、浪、流等扰动对决策过程的影响. 在训练阶段未对扰动进行显式建模, 会导致策略对
环境不确定性的鲁棒性有限, 泛化至真实海况时性能显著下降. 此外, 面向三维复杂环境
(含水面障碍物、近岸浅滩等) 的欠驱动无人艇追逃策略研究相对匮乏, 奖励函数的设计多
停留于距离相关的简单形式, 缺乏对运动姿态与追击过程的精细刻画.

\textbf{问题 2: 多无人艇协同围捕博弈缺乏严格成功判据与实物实验验证.} 现有多无人艇
协同围捕研究普遍存在两方面局限: 一方面, 对 ``围捕成功'' 这一核心评价指标缺乏严格
的数学定义, 研究者多以 ``追击者与逃逸者距离小于阈值'' 作为近似判据, 忽视了包围编
队对逃逸者可达区域的实质约束; 另一方面, 算法设计上缺乏对围捕过程阶段性 (接近、成
型、保持) 的显式刻画, 奖励函数往往采用单一形式, 导致训练过程中出现稀疏奖励、收敛
缓慢等问题. 更重要的是, 绝大多数研究仅完成数值仿真验证, 缺乏面向真实水面环境的多
无人艇物理实验, 算法从仿真到实物的迁移能力和工程可行性尚待考察.

\textbf{问题 3: TAD 博弈中非自杀式多对多防御策略研究空白.} 现有 TAD 博弈研究大多
针对简化的三方场景或自杀式防御者, 对于非自杀式约束下多防御者对多进攻者的博弈策
略研究十分稀缺. 此外, 现有工作通常仅以即时观测作为决策依据, 未能有效利用进攻者历
史状态序列与未来状态预测, 难以适应动态复杂的对抗场景; 在奖励设计上, 对多防御者之
间的协同行为缺乏系统性建模, 策略收敛速度与稳定性均有待提升.

\section{本文主要研究内容与创新点}
\label{sec:contribution}
针对上述现有研究存在的问题, 本文以欠驱动水面无人艇为研究对象, 围绕复杂海洋环境下
的无人艇博弈策略优化展开系统研究, 主要工作与创新点如下:

\textbf{(1) 基于生物启发强化学习的欠驱动无人艇追逃博弈策略.} 针对含三维障碍物与环
境扰动的复杂水域下欠驱动无人艇追逃决策困难的问题, 提出一种基于近端策略优化算法的
生物启发式追逃博弈策略. 主要创新点包括: (i) 设计了融合运动观测、环境感知观测与\textbf{
扰动观测}的三段式观测空间, 将风浪流等外部扰动显式纳入决策依据, 有效提升了策略对环
境不确定性的鲁棒性; (ii) 受鲨鱼捕食行为启发, 提出融合距离逼近、角度对齐与避障惩罚
的复合奖励函数, 缓解了稀疏奖励问题, 加快了训练收敛速度; (iii) 在包含不规则障碍物的
三维仿真环境中完成了追逃博弈训练, 并通过多组扰动条件下的对比实验验证了所提算法的
有效性与鲁棒性. 相关内容详见本文第三章.

\textbf{(2) 基于威胁势场与多头注意力机制的多无人艇协同围捕博弈策略.} 针对协同围捕
博弈中成功判据缺乏、阶段特征难刻画、工程验证不足等问题, 提出一种结合威胁势场与多
头注意力机制的多无人艇协同围捕博弈算法. 主要创新点包括: (i) 基于六自由度非线性运
动学模型, 首次给出了协同围捕成功的显式数学判据, 建立了包围区域与逃逸者可达区域之
间的严格集合包含关系; (ii) 提出能量波峰 (Energy Crest) 与能量波谷 (Energy
Trough) 概念, 将围捕过程划分为接近、成型、保持三阶段, 设计了与阶段自适应相关的
密集奖励机制; (iii) 在多智能体近端策略优化 (MAPPO) 算法的 Critic 网络中引入多头
注意力机制, 构建了 MHD-MAPPO 算法, 显著提升了多艇协同决策效率; (iv) \textbf{搭
建了多无人艇物理实验平台, 在真实水面环境下完成了 2 对 1 与 3 对 1 协同围捕实物实
验}, 从仿真到实物两个层面全面验证了所提算法的工程可行性. 相关内容详见本文第四章.

\textbf{(3) 基于时序感知与预测的非自杀式 TAD 目标防御博弈策略.} 针对现有 TAD 博弈
研究未系统考虑非自杀式约束与进攻者状态预测的问题, 提出一种基于时序感知与预测的
后见经验分配算法 (TP-POCA). 主要创新点包括: (i) 建立了动态目标护航与静态目标固守
两阶段 TAD 任务模型, 系统给出了非自杀式防御者与进攻者的策略约束; (ii) 在策略网络
中引入残差自注意力模块以建模多防御者间的相关性, 并利用循环神经结构捕捉观测的时序
依赖; (iii) 基于扩展卡尔曼滤波 (EKF) 构建进攻者状态预测模块, 使防御者具备前瞻性决
策能力; (iv) 受鲸群协同防御行为启发, 设计了势场化的奖励函数, 显著提升了多防御者
的协同水平; (v) 在动态与静态目标场景下分别开展了大量仿真实验与消融分析, 验证了所
提算法在防御成功率与策略稳定性方面的优势. 相关内容详见本文第五章.

\section{本文组织结构}
\label{sec:organization}
为系统呈现本文的研究工作, 全文共分为六章, 章节安排如下:

\textbf{第一章 \ 绪论.} 介绍本文课题的研究背景与意义, 梳理国内外关于水面无人艇智能
控制、无人艇博弈策略、多智能体深度强化学习以及目标-进攻-防御博弈的研究现状, 凝练
现有研究中存在的关键问题, 明确本文的研究内容、创新点与组织结构.

\textbf{第二章 \ 预备知识与理论基础.} 介绍本文研究所需的基础理论, 包括水面无人艇六
自由度运动学与动力学模型、海洋环境扰动建模、博弈论基础 (微分博弈、马尔可夫博弈、
部分可观马尔可夫决策过程) 及深度强化学习相关算法 (PPO、MAPPO、SAC、注意力机制等),
为后续章节的算法设计提供理论支撑.

\textbf{第三章 \ 基于生物启发强化学习的无人艇追逃博弈策略.} 围绕三维复杂环境下的欠
驱动无人艇追逃博弈问题, 建立了包含扰动观测的三段式观测空间, 提出了鲨鱼捕食行为启
发的复合奖励函数, 基于 PPO 算法构建了追逃博弈策略, 并通过仿真实验验证其在多种扰
动条件与典型场景下的有效性和鲁棒性.

\textbf{第四章 \ 基于威胁势场与多头注意力机制的多无人艇协同围捕博弈策略.} 面向复杂
环境下的多无人艇协同围捕博弈, 给出了围捕成功的显式数学判据, 提出能量波峰与能量波
谷概念及三阶段奖励机制, 构建了嵌入多头注意力的 MHD-MAPPO 算法, 完成了多场景仿真
实验与多无人艇物理实验验证.

\textbf{第五章 \ 基于时序感知与预测的非自杀式 TAD 目标防御博弈策略.} 针对非自杀式
约束下的 TAD 博弈问题, 建立了动态目标护航与静态目标固守的两阶段任务模型, 提出了
融合残差自注意力、循环时序结构、扩展卡尔曼滤波预测与后见经验分配的 TP-POCA 算法,
通过丰富的仿真实验与消融研究验证了算法的有效性.

\textbf{第六章 \ 总结与展望.} 总结全文研究工作, 归纳主要创新点, 并对未来可扩展的研
究方向进行展望.

本文总体研究框架如图~\ref{fig:thesis_framework} 所示.

\begin{figure}[!htbp]
  \centering
  % \includegraphics[scale=0.45]{thesis_framework.pdf}
  \vspace{4cm}
  \bicaption{本文总体研究框架}{Overall research framework of this thesis}
  \label{fig:thesis_framework}
\end{figure}
